\lecture{5}{8. September 2025}{Rigid Body Equilibrium, pt. 1}

\section{Equilibrium of a Rigid Body}

\subsection{Conditions for Rigid-Body Equilibrium}
As mentioned previously, any system of forces and couple moments acting on a body can be reduced to an equivalent resultant force and resultant couple moment about any arbitrary point $O$ on or off the body. If these two resultants are both equal to zero, then the body is said to be in equilibrium. Mathematically this is:
\begin{align*}
  \textbf{F}_R &= \sum \textbf{F} = \textbf{0} \\
  \left( \textbf{M}_R\right)_O &= \sum \textbf{M}_O = \textbf{0}
.\end{align*}
These two equations are both necessary and sufficient for equilibrium. Also, this is true in both 2 and 3 dimensions.


\subsection{Characteristics of Dry Friction}
\textit{Friction} is the name of the force that resists the movement between two contacting surfaces that slide relative to each other. Here, only \textit{dry friction}, also called \textit{Coulomb friction} will be analyzed. Dry friction occurs between two bodies when there is no lubricating fluid. 

The friction force will work to impede motion of an external force $\textbf{P}$. For small values of $\textbf{P}$ the friction force might be big enough to totally impede motion and induce equilibrium. However if $\textbf{P}$ is increased, then the force of friction $\textbf{F}$ will also increase until it attains a maximum value $F_s$ called the \textit{limiting static frictional force}. This is given by:
\[ 
F_s = \mu_s N
\]
where $N$ is the normal force and $\mu_s$ is the coefficient of static friction. 

After sliding occurs the frictional force at the contacting surface will drop to a smaller value $F_k$ called the \textit{kinetic frictional force}. The formula for this is:
\[ 
F_k = \mu_k N
\]
Where $\mu_k$ is the coefficient of kinetic friction.
